% ЗАКЛЮЧЕНИЕ, без номер.
\section*{ЗАКЛЮЧЕНИЕ}
\label{sec:conclusion}
\addcontentsline{toc}{section}{ЗАКЛЮЧЕНИЕ}

Проектът разработва система за анализ на транспортната сигурност и за управление на допускането
на връзки, изградена от отделни слоеве и проверена в изолирана собствена среда. Анализаторът
разчита наблюдаемата част от ръкостискането по TLS 1.3 със скорост 128\,776 обмена в секунда на
едно ядро. Проверката на веригите е реализирана самостоятелно по алгоритъма от \cite{rfc5280} и
дава очаквания изход за всичките тринадесет построени случая. Четирите механизма за допускане се
включват поотделно, което позволи цената им да бъде измерена, а не предположена.

\textbf{Основният измерен резултат.} Собствената цена на едно решение е между 9 и 148
наносекунди, а цената на една допусната връзка под натоварване е между 454 и 1094 микросекунди.
Следователно намалението на пропускателната способност при включена бисквитка, до 44,6 на сто от
базата, не идва от изчислението, а от допълнителния обмен. Срещу тази цена върховото запълване
на таблицата пада шест пъти, а заделената памет за състояние пет пъти, при едно и също
количество лавинен трафик. Мерките, които се обсъждат като скъпи, се оказаха скъпи в брой
обмени, а не в процесорно време.

\textbf{Ограничения.} Проверката на подписа не се извършва, така че системата не открива верига с
правилни структурни полета и подправен подпис; това е отпечатано при всяка проверка. Пасивният
анализ вижда само част от обмена, което е свойство на протокола. Опитите са на една машина и
върху локалния интерфейс, където всички клиенти делят един адрес, така че ограничаването по
източник не може да бъде оценено срещу разпределен източник. Плавното влошаване е проверено от
тестовете, но не е измерено под натоварване. Машината е обща, поради което един ред в
\tabref{tab:decision} е отбелязан като недостатъчно устойчив. И най-същественото: системата не е
подлагана на атака, на фъзинг или на външен одит, тоест нищо тук не бива да се чете като доказана
устойчивост.

\textbf{По-нататъшна работа.} Проверка на подписа зад условна зависимост, така че да се включва
при налична криптографска библиотека и да се отчита като непроведена, когато я няма. Измерване в
среда с различими източници, което е единственият начин ограничаването по източник и плавното
влошаване да бъдат оценени при условията, за които са предназначени. Фъзинг на разчитачите на TLS
и на DER, тъй като те приемат вход от чужда страна. Разширяване към QUIC, където същата ключова
обмяна се пренася в друг транспорт \cite{rfc9001}.
